\begin{topic}[id=safety_vs_security,
             difficulty={Mittel},
             type={Begriffs- und Systemanalyse},
             prerequisites={Grundlagen IT-Sicherheit; Interesse an autonomen oder sicherheitskritischen Systemen},
             practice={optional},
             owner={Dozententeam},
             updated={2026-04-09},
             tags={ IEC 61508, ISO 26262, STRIDE, Hazard, Threat, Echtzeit, Co-Engineering }]{Safety vs. Security in autonomen Systemen}

\TopicDescription{Safety und Security bedingen einander: Security-Prüfungen können Latenz erhöhen; Safety-Mechanismen neue Angriffsflächen schaffen. Das Thema verzahnt Hazard-/Risikobetrachtungen (IEC\,61508/ISO\,26262) mit Threat-Modellen (z.~B. STRIDE) und zeigt Zielkonflikte, Verantwortlichkeiten und Evidenz-Konzepte (Coverage, Fault-Injection, Pen-Tests). Ergebnis ist ein Co-Engineering-Vorgehen mit klaren Artefakten und nachvollziehbaren Entscheidungen.}

\TopicLeitfragen{\item Wie koppelt man Hazard- und Threat-Analysen praktisch?
\item Wo entstehen Zielkonflikte zwischen Echtzeit und Security?
\item Welche Evidenz genügt wofür?
}
\TopicScope{\item Keine Rechts-/Zertifizierungsberatung.
\item Keine tiefen Norm-Exegesen.
}
\TopicPractice{Maßnahmen-Matrix für einen mobilen Roboter (je 3 Safety-/Security-Risiken).}

\TopicKeywords{IEC 61508, ISO 26262, STRIDE, Hazard, Threat, Echtzeit, Co-Engineering}

\nocite{iec61508Overview,iso26262Overview}
\end{topic}
